% Part: proof-theory
% Chapter: natural-deduction
% Section: quantifiers

\documentclass[../../../include/open-logic-section]{subfiles}

\begin{document}

\olfileid[pt]{pt}{nat}{qua}

\olsection{\usetoken{P}{proof} regulares e substituição}

As regras dos quantificadores introduzem algumas complicações devido
às ``condições de variável própria'' que se aplicam às regras
\Elim{\lexists} e~\Intro{\lforall}. A maioria dessas complicações pode
ser resolvida garantindo que a !!{constant}~$c$ nessas inferências
nunca seja reutilizada. Introduzamos a seguinte definição.

\begin{defn}
Uma !!{proof}~$\delta$ em dedução natural é \emph{regular} se
\begin{enumerate}
  \item nenhuma variável própria~$a$ é usada como variável própria de
  mais de uma inferência \Elim{\lexists} ou~\Intro{\lforall}; e
  \item se $a$ é a variável própria de uma inferência \Elim{\lexists}
  ou~\Intro{\lforall}, ela ocorre apenas na sub-!!{proof} imediata que
  conduz, respectivamente, à premissa menor ou à premissa.
\end{enumerate}
\end{defn}

\begin{lem}\ollabel{lem:subst-regular} Se $\delta$ é regular, termina
em~$!C$, $c$ é !!a{constant} não usada como variável própria
em~$\delta$, e $t$~é um termo que não contém nenhuma variável própria
de~$\delta$, então $\Subst{\delta}{t}{c}$, resultante de~$\delta$ pela
substituição de toda ocorrência de~$c$ por~$t$, é uma !!{proof}
regular. A !!{formula} final de $\Subst{\delta}{t}{c}$ é
$\Subst{!C}{t}{c}$. Se $!A^x$ é uma hipótese aberta de~$\delta$, então
$\Subst{!A}{t}{c}^x$ é uma hipótese aberta de~$\Subst{\delta}{t}{c}$.
\end{lem}

\begin{proof}
  Por indução em $\pheight{\delta}$. Se $\pheight{\delta} = 0$, ela
  consiste simplesmente em uma hipótese~$!A^x$. Então,
  $\Subst{\delta}{t}{c}$ é~$\Subst{!A}{t}{c}^x$.

  Se $\pheight{\delta} > 0$, distinguimos casos segundo a última
  inferência em~$\delta$. Os casos das regras para os conectivos
  proposicionais são fáceis e ficam como exercícios. Consideramos os
  casos em que $\delta$ termina em uma inferência de quantificador.
  \begin{enumerate}
    \item Se a última inferência é $\Intro{\lforall}$, $\delta$ tem a
    forma
    \[
    \AxiomC{}
    \RightLabel{$\delta'$}
    \DeduceC{$!A(a,c)$}
    \RightLabel{\Intro{\lforall}}
    \UnaryInfC{$\lforall[x][!A(x,c)]$}
    \DisplayProof
    \]
    Pela hipótese de indução, $\Subst{\delta'}{t}{c}$ é uma
    !!{proof} regular de $!A(a,t)$. Como $a$ é uma variável própria
    de~$\delta$, $a$~não ocorre em~$t$. Portanto, a última inferência em
    $\Subst{\delta}{t}{c}$,
    \[
    \AxiomC{}
    \RightLabel{$\Subst{\delta'}{t}{c}$}
    \DeduceC{$!A(a,t)$}
    \RightLabel{\Intro{\lforall}}
    \UnaryInfC{$\lforall[x][!A(x,t)]$}
    \DisplayProof
    \]
    está correta: como $t$ não contém~$a$, nem a conclusão nem qualquer
    hipótese aberta contém~$a$.
    \item Se a última inferência é $\Elim{\lforall}$,
    $\delta$ tem a forma
    \[
    \AxiomC{}
    \RightLabel{$\delta'$}
    \DeduceC{$\lforall[x][!A(x,c)]$}
    \RightLabel{\Elim{\lforall}}
    \UnaryInfC{$!A(s(c),c)$}
    \DisplayProof
    \]
    Pela hipótese de indução, $\Subst{\delta'}{t}{c}$ é uma
    !!{proof} regular de $\lforall[x][!A(x,t)]$. (O termo~$s(c)$ na
    conclusão pode conter~$c$, mas não necessariamente.) A última
    inferência em
    $\Subst{\delta}{t}{c}$,
    \[
    \AxiomC{}
    \RightLabel{$\Subst{\delta'}{t}{c}$}
    \DeduceC{$\lforall[x][!A(x,t)]$}
    \RightLabel{\Elim{\lforall}}
    \UnaryInfC{$!A(s(t),t)$}
    \DisplayProof
    \]
    está correta, e $!A(s(t),t) \ident \Subst{!A(s(c),c)}{t}{c}$. O
    caso de \Intro{\lexists} é semelhante.
  \item Se a última inferência é $\Elim{\lexists}$, $\delta$ tem a
    forma
    \[
    \AxiomC{}
    \RightLabel{$\delta_1'$}
    \DeduceC{$\lexists[x][!A(x,c)]$}
    \AxiomC{$\Discharge{!A(a,c)}{x}$}
    \RightLabel{$\delta_2'$}
    \DeduceC{$!C$}
    \DischargeRule{\Elim{\lexists}}{x}
    \BinaryInfC{$!C$}
    \DisplayProof
    \]
    Pela hipótese de indução, $\Subst{\delta_1'}{t}{c}$ e
    $\Subst{\delta_2'}{t}{c}$ são !!{proof}s regulares de
    $\lexists[x][!A(x,c)]$ e de $!C$ a partir da hipótese
    aberta~$!A(a,t)$. Como $a$ é uma variável própria de~$\delta$,
    $a$~não ocorre em~$t$. Portanto, a última inferência em
    $\Subst{\delta}{t}{c}$,
    \[
    \AxiomC{}
    \RightLabel{$\Subst{\delta'}{t}{c}$}
    \DeduceC{$!A(a,t)$}
    \RightLabel{\Intro{\lforall}}
    \UnaryInfC{$\lforall[x][!A(x,t)]$}
    \DisplayProof
    \]
    está correta: $\Subst{!A(x,t)}{a}{x} \ident
    \Subst{!A(a,c)}{t}{c}$; nem a conclusão~$!C$ nem qualquer hipótese
    aberta contém~$a$.
  \end{enumerate}
\end{proof}

\begin{prop}\ollabel{prop:regular}
Se $\delta$ é !!a{proof} em \Log{N1c} ou \Log{N1i}, então existe uma
!!{proof} regular~$\delta'$ com a mesma estrutura (em particular, a
mesma !!{height}) que difere de $\delta$ apenas pela substituição de
variáveis próprias.
\end{prop}

\begin{proof}
Somente para os fins da demonstração, chame uma inferência de variável
própria de~$\delta$ de \emph{limpa} se sua variável própria não é a
variável própria de outra inferência nem ocorre em qualquer lugar
de~$\delta$, exceto na sub-!!{proof} imediata da inferência; caso
contrário, chame-a de \emph{suja}. Seja $n$ o número de inferências de
variável própria sujas em~$\delta$. Mostramos, por indução em~$n$, que
$\delta$~pode ser transformada na !!{proof} regular exigida~$\delta'$.
Se $n=0$, todas as inferências de variável própria em~$\delta$ são
limpas; portanto, $\delta$ é regular e não há nada a demonstrar. Caso
contrário, escolha uma inferência de variável própria mais alta,
digamos \Intro{\lforall}. A sub-!!{proof}~$\delta_1$ de~$\delta$ que
termina nela tem a forma
    \[
    \AxiomC{}
    \RightLabel{$\delta_2$}
    \DeduceC{$!A(c)$}
    \RightLabel{\Intro{\lforall}}
    \UnaryInfC{$\lforall[x][!A(x)]$}
    \DisplayProof
    \]
Observe que, como $c$ é uma variável própria, ela não ocorre
em~$\lforall[x][!A(x)]$ nem em uma hipótese aberta. A prova $\delta_2$
é regular, pois não contém nenhuma inferência de variável própria.
Seja $a$ !!a{constant} que não ocorre em~$\delta$. Por
\olref{lem:subst-regular}, $\Subst{\delta_2}{a}{c}$ é uma prova regular
de $!A(a)$. Pela condição de variável própria, nenhuma hipótese aberta
de~$\delta_2$ contém~$c$; portanto, nenhuma hipótese aberta de
$\Subst{\delta_2}{a}{c}$ contém~$a$. Assim, podemos
aplicar~$\Intro{\lforall}$. Se substituirmos $\delta_1$ em~$\delta$
pela prova $\delta_1'$,
    \[
    \AxiomC{}
    \RightLabel{$\Subst{\delta_2}{a}{c}$}
    \DeduceC{$!A(a)$}
    \RightLabel{\Intro{\lforall}}
    \UnaryInfC{$\lforall[x][!A(x)]$}
    \DisplayProof
    \]
teremos substituído uma inferência de variável própria suja por uma
limpa, e a única diferença é a substituição da variável própria $c$
por~$a$. A prova resultante tem $n-1$ inferências de variável própria
sujas, e o resultado segue da hipótese de indução.
\end{proof}

\begin{prob}
Descreva o caso da demonstração de \olref{prop:regular} em que a
inferência de variável própria mais alta é~\Elim{\lexists}.
\end{prob}

Não recorreremos explicitamente à proposição que acabamos de
demonstrar, mas assumiremos que todas as !!{proof}s consideradas são
regulares---se não forem, sempre poderemos substituir as variáveis
próprias para torná-las regulares.

\Olref{lem:subst-regular} e \olref{prop:regular} também valem para
\Log{N2c} e~\Log{N2i}. Deixamos as demonstrações como exercícios.

\end{document}
